% Exported from Math++. Compile with pdflatex.
\documentclass[11pt,a4paper]{article}
\usepackage[T1]{fontenc}
\usepackage[utf8]{inputenc}
\usepackage{lmodern}
\usepackage{amsmath,amssymb,amsthm,mathtools}
\usepackage{graphicx,booktabs,xcolor}
\usepackage[hidelinks]{hyperref}
\usepackage[margin=25mm]{geometry}
\theoremstyle{plain}
\newtheorem{theorem}{Theorem}[section]
\newtheorem{lemma}[theorem]{Lemma}
\newtheorem{corollary}[theorem]{Corollary}
\newtheorem{proposition}[theorem]{Proposition}
\newtheorem{conjecture}[theorem]{Conjecture}
\theoremstyle{definition}
\newtheorem{definition}[theorem]{Definition}
\newtheorem{example}[theorem]{Example}
\theoremstyle{remark}
\newtheorem*{remark}{Remark}
\setlength{\parindent}{0pt}\setlength{\parskip}{0.6em}
\begin{document}
\title{Projectile motion on a slope}
\author{Ben Elmore}
\date{20 November 2026}
\maketitle
\begin{abstract}
We find how far a projectile launched at speed $v_0$ and angle $\theta$ travels, following the classical treatment \cite{halliday14} and the numerical study of \cite{lee20,knuth84}.
\end{abstract}
\section{Introduction}

Every worksheet works without any of this: the title block, numbered theorems, citations and notes are all optional.

\section{Range on flat ground}

Let $g \equiv 9.81$ and $v_0 \equiv 20$. The range is given in \eqref{range}.

\begin{equation}
  R(\theta) \coloneqq  \frac{v_0^2 \sin 2\theta}{g} \label{range}
\end{equation}
\begin{equation*}
  R(\pi / 4) = \frac{40000}{981}
\end{equation*}
\begin{definition}
The \emph{range} of a launch is the horizontal distance to the point where the projectile lands again.

\end{definition}
\begin{theorem}
For a fixed launch speed the range is largest when $\theta = \pi/4$.

\end{theorem}
\begin{proof}
By \eqref{range} the range is proportional to $\sin 2\theta$, which is at most $1$ and equals $1$ exactly when $2\theta = \pi/2$.\footnote{Strictly, for $0 \le \theta \le \pi/2$.}
\end{proof}
\begin{remark}
Air resistance is ignored throughout.

\end{remark}
\section{Slopes}

\begin{lemma}
A slope shortens the range of an uphill launch.

\end{lemma}
\begin{corollary}
The best angle on a slope is below $\pi/4$.

\end{corollary}
\begin{thebibliography}{99}
\bibitem{knuth84} Knuth, D. E. \emph{The TeXbook}, 1984. \url{https://example.org/texbook}
\bibitem{lee20} Lee, K. and Park, H. \emph{Slopes and projectiles}, 2020.
\bibitem{halliday14} Halliday, D.; Resnick, R.; Walker, J. \emph{Fundamentals of Physics}, 2014.
\end{thebibliography}
\end{document}
